\addcontentsline{toc}{chapter}{TÓM TẮT}
\vspace*{0.5cm}
\begin{center}
    \textbf{TÓM TẮT}
\end{center}

Điều khiển cánh tay robot từ xa thông qua thực tế ảo (VR Teleoperation) là một hướng nghiên cứu đang được quan tâm mạnh mẽ, khi nhu cầu vận hành thiết bị từ xa và an toàn ngày càng trở nên cấp thiết. Mục tiêu của việc xây dựng một ứng dụng từ xa hiện đại với VR, đó là tăng cường hiệu quả của tương tác cộng tác Người - Máy (Human-Robot Collaboration), đem lại trải nghiệm điều khiển dễ dàng, ổn định và thật sự có ý nghĩa thực tiễn cho người dùng

Dựa trên các kết quả nghiên cứu hiện nay, đề án nỗ lực để xây dựng một hệ thống điều khiển cánh tay robot từ xa thông qua kính VR, với mục tiêu sử dụng những công nghệ mới nhất, thiết kế và tích hợp chúng với nhau một cách hoàn chỉnh. Chủ đề: "Nghiên cứu và phát triển hệ thống điều khiển cánh tay robot với ROS2 đồng bộ kính thực tế ảo", hướng tới việc xây dựng một nền tảng đầy đủ từ cấp thấp lên cao, bao gồm xử lý việc giao tiếp với thiết bị cánh tay robot công nghiệp cũ - VS-6577 cùng bộ điều khiển RC5 giới hạn về cách thức giao tiếp, xây dựng hệ thống điều khiển sử dụng ROS2 tích hợp với Ros2-Control và MoveIt2, tuân thủ chặt chẽ theo kiến trúc ứng dụng được đề xuất bởi nhà phát triển. Từ một nền tảng điều khiển vững chắc, ứng dụng VR xây dựng Digital Twin của cánh tay robot được phát triển bằng Unity, chạy trên kính Meta Quest 3 một cách độc lập, không phụ thuộc vào máy tính, với mục tiêu đem lại trải nghiệm điều khiển từ xa hiệu quả, trực quan và ổn định thông qua kính VR.

Về cơ bản, đề án đã thực hiện và hoàn thành mục tiêu phát triển và thiết kế đã đề ra, tất cả các module đều hoàn tất và có những đánh giá phù hợp. Ứng dụng chạy ổn định với khả năng điều khiển tốt, phù hợp cho khả năng sử dụng và vận hành lâu dài.

\underline{Các từ khóa:} ROS 2, Ros2-Control, MoveIt2, Điều khiển từ xa, Thực tế ảo, Meta Quest 3, Denso VS-6577, RosBridge, Digital Twin, Ros-Sharp.

\newpage
\addcontentsline{toc}{chapter}{ABSTRACT}
\vspace*{0.5cm}

\begin{center}
    \textbf{ABSTRACT}
\end{center}

Remote robotic arm control through virtual reality (VR teleoperation) has become a major research direction, as the need for safe and reliable remote operation is increasingly urgent. The goal of developing a modern VR-based remote-control application is to improve Human-Robot Collaboration (HRC) effectiveness and provide an intuitive, stable, and practically meaningful control experience for users.

Inspired by many VR-Robot related works, this project strives to build a VR-based teleoperation system for a robotic arm by adopting state-of-the-art technologies and integrating them into a complete architecture. The topic, "Research and Development of a ROS2-Based Robotic Arm Control System Synchronized with a VR Headset," aims to establish a full-stack platform from low-level to high-level layers. The work includes communication with a legacy industrial robotic arm (Denso VS-6577) and an RC5 controller with constrained communication methods, and the development of a ROS2 control stack integrated with Ros2-Control and MoveIt2, in close alignment with the software architecture recommended by the Ros2-Control developers. Built on this solid control foundation, a VR application implementing a Digital Twin of the robotic arm was developed in Unity and deployed to run natively on Meta Quest 3 without dependence on a PC, with the goal of delivering effective, intuitive, and stable teleoperation through a VR headset.

In summary, the thesis has achieved the proposed design and development objectives. All modules were completed and evaluated with appropriate criteria. The resulting application operates stably, provides good control performance, and is suitable for long-term use and operation.

\underline{Keywords:} ROS 2, Ros2-Control, MoveIt2, Teleoperation, Virtual Reality, Meta Quest 3, Denso VS-6577, RosBridge, Digital Twin, ROS-Sharp.
